\chapter{Jets as a probe of a Quark-Gluon Plasma\label{PreviousCorrelationResults}}% needs a better title
%Thorsten Renk's paper http://arxiv.org/abs/hep-ph/0608333 has the pretty plots in B&W
%the color figures are
%http://fpaxp1.usc.es/nestor/session6/Renk.pdf

%Need to add Wenger's results

Due to confinement, particles which carry a color charge, such as quarks and gluons, cannot be directly observed.  Instead, they fragment into hadrons, which are color-neutral, and these hadrons are detected.  In hard scattering processes the majority of the hadrons are produced roughly colinear to the fragmenting partons, resulting in a narrow cone of hadrons called a jet.  At RHIC energies most jets are produced in 2$\rightarrow$2 scattering of hard partons from the incoming nuclei.  Since the incoming partons have momenta nearly parallel with the nuclei, momentum conservation causes the scattered partons to be separated by roughly 180\textdegree\xspace in azimuth.  Events with jets at RHIC energies are dominantly back-to-back di-jets.  \Fref{ppJetSample} shows an example of a di-jet event in STAR.  The beam goes through the center and the lines show the different sectors of the Time Projection Chamber.

\begin{figure}
  \centering\resizebox{8cm}{!}{
  \includegraphics{/home/madscientistgirl/HarrisLab/Thesis/latex/previousresults/ppjet.eps}}\\

  \caption[Sample di-jet in a \pp collision at \sNNtwohundred]
  {Sample jet observed in STAR in a real \pp collision at \sNNtwohundred
  \label{ppJetSample}}
\end{figure}


The jet cross section can be calculated in perturbative QCD as the average of quark, anti-quark, and gluon processes.  For two partons scattering the jet cross section is given by
\begin{equation}
 \sigma_{i,j\rightarrow k} = \sum_{i,j}\int dx_1 dx_2 dt f^1_i(x_1,Q^2) f^2_j(x_2,Q^2) \frac{d\sigma_{i,j\rightarrow k}}{dt}\label{jetcrosssection}
\end{equation}
where $x_1$ and $x_2$ are the fraction of the nucleon's momentum carried by the parton, $Q^2$ is the momentum transfer in the scattering, $f^1_i$ and $f^2_j$ are the parton distribution functions in the nucleons, and $\sigma_{i,j\rightarrow k}$ is the cross section for the reaction $ij\rightarrow k$ \cite{Peskin:257493}.  Parton distribution functions are the probability densities for finding a parton as a function of $Q^2$ and x.

After the hard scattering, the scattered partons hadronize.  The process of hadronization is non-perturbative.  It is described by fragmentation functions, which represent the probability for a parton to fragment into a particular hadron.  Fragmentation functions have not been determined from first principles, however, it is possible to parameterize them as functions of the parton energy and the fraction of energy carried by the parton and then determine the fragmentation functions from data.  Vacuum fragmentation functions have been well constrained by data from elementary collisions.


\begin{figure}
  \centering\resizebox{7cm}{!}{
  \includegraphics{/home/madscientistgirl/HarrisLab/Thesis/latex/previousresults/idjetschematic.eps}}\\

  \caption[Schematic diagram of a di-jet]
  {Schematic diagram of a di-jet in a heavy ion collision.  The red arrow denotes the leading hadron, the one with the highest \pT, and the black arrows denote other hadrons produced during hadronization.
  \label{IDJetSchematic}}
\end{figure}


In heavy ion collisions, jets are formed from parton scattering which likely occurs early in the collision.  The partons then travel through the medium and therefore serve as a probe of the medium.  This scenario is depicted in \Fref{IDJetSchematic}.   Since jets have been studied extensively in \pp collisions, modifications of jets in \AplusA collisions relative to \pp can be attributed to interactions with the medium.  These modifications may come in the form of parton energy loss before fragmentation or modification of the fragmentation functions for partons that fragment in the medium.

There have been many different approaches to studying jets at RHIC.  In this chapter the various methods and key results from RHIC are reviewed.  Some of these studies occurred concurrently with the analyses presented in this thesis, and discussion of these results is necessary to form a complete picture of what can be learned from the data on jets at RHIC.

\section{Studies of jets through high-\pT triggered di-hadron correlations\label{Correlations}}

%Full jet reconstruction in a heavy ion environment is difficult and while studies have recently been done \cite{jetrecoSevilWWND,jetrecoJoernHP,jetrecoSevilHP,jetrecoSevilHQ}, initial studies of jets were done by looking at correlations between high-\pT hadrons.  These studies remain valuable even now that full jet reconstruction has been done in heavy ion collisions because the biases are different in the two different methods.  Full jet reconstruction may be more biased towards unmodified jets because jet finding algorithms are dependent on jet topologies similar to those in \ee and \pp collisions.  In addition, studies that may be limited by statistics, such as studies of multistrange particles in jets, may not be feasible using jet reconstruction.

In di-hadron correlations, a high momentum particle is selected and the distribution of particles relative to that particle is determined.  The former particle is called the trigger particle and the latter are called the associated particles.  The primary criterion used to determine trigger and associated particles is their momenta and the method neglects contributions from any sources of correlations between high-\pT particles other than jets and anisotropic flow.  Anisotropic flow is a background and its subtraction is discussed further in \Cref{CorrMethod}.  The momenta of the trigger and associated particles are restricted to high-\pT to increase the probability that the particles come from a jet and therefore decrease the combinatorial background.  It has been argued that contributions from other sources are non-negligible.  In particular it has been proposed that a significant fraction of trigger particles may not come from jets \cite{Jia:2008kf}.  The STAR collaboration presented correlations normalized per trigger particle and with this normalization, results from different systems (\pp, \dAu, \Cu, and \Au) would be identical if there were no modification of the jet\cite{STARMiniJets2006,STARPunchThrough,STARLowPtAssoc2005,STARRxnPlane2004,STARIconicJetQuench2003,STARIconicJetQuench2002,STARJetQuench2002}.  The PHENIX collaboration typically used a different normalization where the amplitude of the correlation is interpreted as the probability for an associated particle to be correlated with the jet \cite{PHENIXPartIDCorr2007,PHENIXAuAu200Y4Corr2008,PHENIXDandruff,PHENIXCorrEnergyDependence,PHENIXppCorr,PHENIXAuAu200AwaySideCorr,PHENIXAuAu200Y2Corr}.


\subsection{Early studies of jets at RHIC}


\begin{figure}
  \centering\resizebox{10cm}{!}{
  \includegraphics{/home/madscientistgirl/HarrisLab/Thesis/latex/previousresults/IconicJetQuench.eps}}\\

  \caption[Jet quenching observed in di-hadron correlations]
  {Di-hadron correlations from \pp, \dAu, and \Au collisions  at \sNNtwohundred \cite{STARIconicJetQuench2003}.  (a)  Correlations for \pp, minimum bias \dAu, and central \dAu collisions with no background subtraction. (b) Correlations for \pp, \dAu, and 0-10\% central \Au collisions after background subtraction.  Data are for 2 \GeV $<$ \ptassoc $<$ 4 \GeV and 4 \GeV $<$ \pttrig $<$ 6 \GeV.
  \label{IconicJetQuench}}
\end{figure}

\Fref{IconicJetQuench} shows di-hadron correlations from STAR  \cite{STARIconicJetQuench2003}.  The peak near \dphi = 0 comes from hadrons which fragmented from the same parton as the trigger hadron and is called the \ns jet.  The peak near \dphi \roughly $\pi$ comes from a momentum conserving jet opposite the \ns jet and is called the \as jet.  \Fref{IconicJetQuench}(a) compares correlations from \pp, minimum bias \dAu, and the 0-20\% highest multiplicity \dAu collisions and demonstrates that while the background is higher in \dAu, there are no dramatic differences between \pp and \dAu.  This implies that initial state effects do not significantly modify the shape of jets, which would be evident in the \dAu data. Even though there is no jet quenching in \dAu, the \as is broadened in pseudorapidity.  This effect is present in simple models of \pp collisions and is due to (1) the difference between the center of mass frame of the scattered partons and the lab frame and (2) the difference between the jet axis and the direction of the momentum of the leading parton.  \Fref{IconicJetQuench}(b) compares \pp, \dAu, and central \Au after background subtraction.  Whereas there is a clear \ns peak in all systems, the \as peak is not evident in central \Au data.  The disappearance of the \as jet is called jet quenching.  High-\pT partons are biased towards the surface, as shown in \Fref{SurfaceBias}.  The surface bias is expected to be greater for di-hadron correlations because two high-\pT partons are required \cite{Renk:2006pw}.  This is significant for the interpretation of di-hadron correlations, as it implies that the \ns jet is more likely to come from the surface.

Theoretical studies of correlations usually treat the \ns and the \as separately.
It is generally assumed that the \ns comes from surface-biased partons, while the \as comes from partons which have traveled through the medium.

\subsection{The Away-side}

\begin{figure}
  \centering\resizebox{15cm}{!}{
  \includegraphics{/home/madscientistgirl/HarrisLab/Thesis/latex/previousresults/LowPtCorrelations.eps}}\\

  \caption[Low-\pT correlations]
  {Di-hadron correlations after background subtraction for 0.15 \GeV $<$ \ptassoc $<$ 4 \GeV and 4 \GeV $<$ \pttrig $<$ 6 \GeV from \Au  collisions at \sNNtwohundred \cite{Horner:2006yd}.
  \label{LowPtCorr}}
\end{figure}


\Fref{IconicJetQuench} shows the di-hadron correlation for relatively high-\pT associated particles, however, the momentum contained  in the jet must be conserved so there should be some evidence that the \as jet was produced even if it is quenched.  \Fref{LowPtCorr} shows  di-hadron correlations for lower momentum associated particles.  The \as reappears at lower \pT, but it is broader than in \pp or \dAu collisions (not shown, \cite{Horner:2006yd}).  Many mechanisms were proposed for the production of this structure, including global momentum conservation \cite{Borghini:2006yk}, \as jets deflected 180\textdegree\xspace  away from the \ns jet \cite{Armesto:2004pt}; \v{C}erenkov radiation of gluons \cite{Koch:2005sx}; and the formation of a  Mach cone as the \as jet moves through the medium \cite{MachOne,CasalderreySolana:2005rf}.  In these models particles on the \as will have different correlations with each other, so evidence for these mechanisms was investigated by  looking at three-particle correlations \cite{Ulery:2008rw,Ulery:2008pj,Ulery:2006ha,Pruneau:2008ip,Pruneau:2007mr}.  3-particle correlations are very sensitive to assumptions about the shape of the background and are statistically limited, therefore, the data are not conclusive as of yet.  However, a simple jet deflection model is ruled out because it cannot explain all of the structures in the data.



\begin{figure}
  \centering\resizebox{10cm}{!}{
  \includegraphics{/home/madscientistgirl/HarrisLab/Thesis/latex/previousresults/PunchThrough.eps}}\\

  \caption[The punch-through phenomenon]
  {Di-hadron correlation before background subtraction with various kinematic cuts from \Au  collisions at \sNNtwohundred \cite{STARPunchThrough}.
  \label{PunchThrough}}
\end{figure}


\fref{PunchThrough} shows that for higher \ptassoc and \pttrig the \as reappears.  At sufficiently high \pT the shape is not dramatically modified relative to the \as in \pp and \dAu collisions, however, the magnitude of the correlation is significantly lower \cite{STARPunchThrough}.  This phenomenon is called punch-through.  It is believed that the selection of high momenta hadrons biases the measurement towards jets produced by partons which have not lost much energy.  This may be because the partons originated from collisions near the surface of the medium or that the initial hard scattering gave the partons enough energy that they could ``punch through'' the medium.


\subsection{The \ns}


\begin{figure}
  \centering\resizebox{10cm}{!}{
\includegraphics{/home/madscientistgirl/HarrisLab/Thesis/latex/previousresults/2DdAu.eps}
}\\
 \centering\resizebox{10cm}{!}{
  \includegraphics{/home/madscientistgirl/HarrisLab/Thesis/latex/previousresults/2DAuAu.eps}
}\\

  \caption[The \ridge in \Au collisions at \sNNtwohundred]
  {Di-hadron correlations in \dAu (upper panel) and 0-12\% central \Au (lower panel) collisions at \sNNtwohundred for 4.0 \GeV $<$ \pttrig $<$ 6.0 \GeV and 2.0 $<$ \ptassoc $<$ \pttrig.  In \dAu there is one peak on the \ns that is narrow in \dphi and \deta called the \jlc.  This peak is evident in the \Au but there is an additional correlation that is narrow in \dphi but broad in \deta called the \ridge \cite{JoernQM06}.
  \label{IconicRidgePlot}}
\end{figure}

The \as was expected to have modifications due to interaction with the medium, however, the \ns is also modified relative to \pp and \dAu collisions.  \Fref{IconicRidgePlot} shows di-hadron correlations as a function of \dphi and \deta in \dAu and \Au collisions.

The \ns peak in \dAu is narrow in both azimuth and pseudorapidity, as expected for vacuum fragmentation.  The peak that is present in \dAu is called the \jlc and is also evident in \Au.  However there is a new component that is narrow in azimuth but broad in pseudorapidity; it is independent of pseudorapidity within the acceptance of the STAR TPC.  This long range pseudorapidity correlation is called the \ridge.  The \ridge has also been observed by PHENIX \cite{McCumber:2008id} and PHOBOS \cite{Wenger:2008ts}.  The measurement in \cite{Wenger:2008ts} demonstrated that the \ridge extends to \deta = 4 for \trigrangevar{2.5} and \assocrangevartwo{0.35}.
%The \ridge has also been observed by PHENIX \cite{McCumber:2008id}, however, these studies are for \assocrange{1.0}{2.0} and \trigrange{2.0}{3.0}.  Within this kinematic region, the \jlc in \Au collisions is considerably wider than in \pp collisions \cite{Bombara:2008zza}, preventing separation of the \jlc and the \ridge.

%\subsection{Correlation Measurements}
\begin{figure}
  \centering\resizebox{9cm}{!}{
  \includegraphics{/home/madscientistgirl/HarrisLab/Thesis/latex/discussion/RidgeVSJetVSNpart.eps}
  }\\

  \caption[The \ridge as a function of \npart.]
  {The number of particles in the \ridge and the \jlc (J+R) and in the \jlc alone (J) as a function of \npart for \trigrange{4.0}{6.0} and \assocrange{2.0}{4.0} within \detarange{-1.75}{1.75} for 0-12\% central \Au collisions at \sNNtwohundred \cite{JoernQM06}.  The solid lines show the systematic error.
%The number of particles in the \jlc is determined using two methods: (1) by subtracting the contribution from the \ridge and (2) determining the yield in azimuth and determining the number of particles in the peak above the plateau in pseudorapidity.  These methods are discussed futher in \Cref{CorrMethod}.
  \label{RidgeVSJetVSNpart}}
\end{figure}

Experimentally, a reasonably complete picture of the \ridge in \Au collisions at \sNNtwohundred can be drawn since the multiplicity in \Au collisions is high and STAR collected a data set with high statistics for \Au collisions at \sNNtwohundred in 2004.  \Fref{RidgeVSJetVSNpart} shows the number of particles in both the \ridge and the \jlc combined and in the \jlc alone as a function of \npart.  The number of particles in the \jlc is determined using two methods: (1) by subtracting the contribution from the \ridge and (2) determining the yield in azimuth and determining the number of particles in the peak above the plateau in pseudorapidity.  These methods are discussed futher in \Cref{CorrMethod}.    The \jlc alone is almost independent of \npart, while the \ridge grows with \npart.  The \ridge is not present in \dAu \cite{JanaQM06}.

The \jly normalized per trigger should remain constant as a function of \npart if jet production in \AplusA collisions is identical to jet production in \pp collisions.  If the yield per trigger changes as a function of \npart this may indicate that jet production is modified by the medium.  \Fref{RidgeVSJetVSNpart} indicates that the \jlc may result from processes that occur in \pp collisions while the \ridge results from processes that either do not occur in \pp collisions or those which are more prominent in \AplusA collisions.

\begin{figure}
  \centering\resizebox{9cm}{!}{
  \includegraphics{/home/madscientistgirl/HarrisLab/Thesis/latex/discussion/RidgeVSPTTrigVSCent2.eps}}\\

  \caption[The number of particles in the \ridge as a function of \pttrig]
  {Number of particles in the \ridge as a function of \pttrig for \assocrange{2.0}{4.0} in \Au collisions at \sNNtwohundred \cite{JoernQM06}.  The solid lines correspond to the systematic errors for the number of particles in the 0-10\% most central collisions, the dashed lines for 30-40\%, and the dotted lines for 40-60\%.
  \label{RidgeVSPTTrigVSCent}}
\end{figure}

\Fref{RidgeVSPTTrigVSCent} shows the number of particles in the \ridge as a function of \pttrig.  The yield of the \ridge is roughly independent of \pttrig and is present even for very high \pttrig.  The number of particles in the \jlc, by comparison, grows with the trigger momentum (not shown, \cite{JanaQM06}), which would be expected for vacuum fragmentation since a higher momentum trigger is more likely to come from a higher momentum jet and therefore the total number of particles produced within the \ptassoc range should be larger.

\begin{figure}
  \centering\resizebox{9cm}{!}{
  \includegraphics{/home/madscientistgirl/HarrisLab/Thesis/latex/previousresults/RidgeVSJetVSPtAssoc.eps}}\\

  \caption[Spectra of particles in the \ridge and the \jlc]
  {Spectra of particles in the \ridge (closed symbols) and the \jlc (open symbols) in 0-12\% \Au collisions at \sNNtwohundred\cite{JoernQM06}.
  \label{RidgeVSJetVSPtAssoc}}
\end{figure}

\Fref{RidgeVSJetVSPtAssoc} shows the spectra of particles in the \ridge as compared to particles in the \jlc and the spectrum of all unidentified hadrons.  The spectrum of particles in the \jlc gets harder the higher the momentum of the trigger particle.  Jets in \pp would lead to the same effect, since higher momentum hadrons are likely to come from higher energy jets.  The \ridge has a spectrum comparable to inclusive unidentified hadrons.  This suggests that the particles in the \ridge may originate from the medium.






\begin{figure}
  \centering\resizebox{12cm}{!}{
  \includegraphics{/home/madscientistgirl/HarrisLab/Thesis/latex/discussion/zTRidgeSubtracted.eps}}\\

  \caption[Fragmentation functions with and without \ridge]
  {Di-hadron fragmentation function \zT = $\frac{p_T}{p_T^{trigger}}$ with and without \ridge subtraction in 0-12\% central \Au collisions at \sNNtwohundred \cite{JoernQM06} (a) shows the di-hadron fragmentation for the \jlc and the \ridge combined.  (b) shows the di-hadron fragmentation function for the \jlc alone after subtracting the \ridge.  (c) shows the ratio of the di-hadron fragmentation function for the \jlc alone to the di-hadron fragmentation function measured in \dAu.
  \label{zTRidgeSubtracted}}
\end{figure}

Traditional studies of fragmentation functions cannot be done precisely without full jet reconstruction because the total parton energy is not known, however, an approximation of the fragmentation function can be made using \pttrig in place of the jet energy \cite{Wang:2003mm}.  This is shown in \Fref{zTRidgeSubtracted} for various trigger momenta.  When the \ridge and the \jlc are combined, there is a clear dependence on \pttrig, however, when the \ridge is subtracted, there is no dependence on \pttrig for the \jlc alone.  Again, this supports a picture where the production of the \jlc is similar to processes that occur in \pp collisions while the \ridge occurs by a novel mechanism.

%ADD JANA's STUFF!!!
\begin{figure}
\begin{tabular}{cc}
   \centering\resizebox{!}{6cm}{
  \includegraphics{/home/madscientistgirl/HarrisLab/Thesis/latex/previousresults/baryon-meson-ratio-inclusive-ridge-jet.eps}} &
   \centering\resizebox{!}{6cm}{
   \includegraphics{/home/madscientistgirl/HarrisLab/Thesis/latex/previousresults/pToPiRatio.eps}} 
 \end{tabular}

  \caption[Particle composition of the \jlc and the \ridge]
  {Comparisons of the \lam/\kaon ratio (left) \cite{JanaWWND2007} and of the \proton/\pion ratio (right) \cite{SuarezQM08} in the \ridge and the \jlc to the inclusive ratios.  Systematic errors are added in quadrature to the statistical errors for the \lam/\kaon ratio and shown as bands for the \proton/\pion ratio.
  \label{RidgeParticleComposition}}
\end{figure}

The composition of the \ridge and the \jlc can also be determined by identifying the associated particles.  \Fref{RidgeParticleComposition} shows the baryon/meson ratios for strange and non-strange particles in the \ridge and \jlc as compared to the inclusive ratios in \Au and \pp as a function of \pT.  The particle ratios in the \jlc are comparable to those in \pp, consistent with a \jlc produced dominantly by fragmentation.  The particle ratios in the \ridge are consistent with the inclusive particle ratios in \Au collisions.  Systematic errors are large, but these data hint at the same trends shown in \Fref{RidgeVSJetVSPtAssoc} - that the \jlc is dominantly produced by fragmentation and that the \ridge comprises particles from the bulk.

No trigger particle type dependence is seen in \Au collisions for either the \ridge or the \jlc \cite{JanaQM06}.  There is no apparent difference between baryon and meson triggers, even though there are dramatic differences between baryons and mesons in nuclear modification factors (shown in \Fref{RAAScaling}) and azimuthal anisotropies (shown in \Fref{PHENIXv2Scaling}).  A trigger particle dependence might be expected if there were a dramatic difference between quark and gluon jets, as the rates of production of \lam and \alam are comparable in gluon jets but many more \lam than \alam are predicted to be produced by quark jets.  No difference between \lam and \alam trigger particles is observed (not shown).  This independence of the yield on the trigger particle type extends even to the \omegabary \cite{Abelev:2007zp}.


%\subsection{Reaction plane dependence}
%This is Aoqi's stuff and I think there's really one figure I need to but I'll patch in two here
\begin{figure}
  \centering\resizebox{15cm}{!}{
  \includegraphics{/home/madscientistgirl/HarrisLab/Thesis/latex/discussion/RxnPlaneDependence0-5Central.eps}
  \includegraphics{/home/madscientistgirl/HarrisLab/Thesis/latex/discussion/RxnPlaneDependence20-60Central.eps}}\\

  \caption[Reaction plane dependence of the \jlc and the \ridge]
  {Reaction plane dependence of number of particles in the \jlc and the \ridge for 0-5\% central (left) and 20-60\% central (right) \Au collisions at \sNNtwohundred \cite{FengQM08}.  The lines correspond to the systematic errors on the number of particles in the \ridge.  The arrow corresponds to the reaction plane averaged \jly in \dAu.  The lines correspond to the systematic errors.
  \label{RxnPlaneDependence}}
\end{figure}

\Fref{RxnPlaneDependence} shows the reaction plane dependence of the \ridge and the \jlcs in two different centrality bins.  The \jlc shows little dependence on the direction relative to the reaction plane, as would be expected if the \jlc originates from vacuum fragmentation.  The \ridge is largest in plane, where the mean path length traveled by the parton would be shortest.  If the \ridge is produced by parton energy loss, these results would indicate that jets produced out of plane are either produced on the surface or completely quenched.

\begin{figure}
  \centering\resizebox{15cm}{!}{
  \includegraphics{/home/madscientistgirl/HarrisLab/Thesis/latex/previousresults/3particleDEtaCorrelations.eps}}\\

  \caption[3-particle correlations on the \ns]
  {3-particle correlations for \trigrange{3.0}{10.0} and \assocrange{1.0}{3.0} for both associated particles on the \ns in \dAu, peripheral \Au, and central \Au collisions at \sNNtwohundred \cite{Netrakanti:2008jw}.
  \label{3partCorrDEta}}
\end{figure}

Additionally, measurements were done to determine if particles in the \ridge are correlated with each other.  3-particle correlations on the \ns are shown in \Fref{3partCorrDEta}.  A trigger particle with \trigrange{3.0}{10.0} is selected and defines the origin.  Only events with at least two associated particles in the range \assocrange{1.0}{3.0} are selected and \Fref{3partCorrDEta} shows the location of the first associated particle relative to the trigger versus the location of the second associated particle relative to the trigger.  A narrow peak around the trigger particle is observed in the \dAu, peripheral \Au, and central \Au data and comes from the same process that produces the \jlc in di-hadron correlations.  In the \Au, this peak sits on top of a plateau; this is the \ridge.  If the two associated particles came from the same source - for instance, from gluon radiation - they would be expected to be correlated in space.  Instead, within the the acceptance of the detector and the statistical errors, particles in the \ridge do not appear to be correlated.  However, an understanding of the effects of the kinematic cuts is required to interpret the results properly.  The lower cut of 1.0 \GeV on \ptassoc is necessary to reduce the combinatorial background, which may make the results more difficult to interpret.  If, for instance, the \ridge were caused by gluon radiation, requiring two particles above 1.0 \GeV could select only those events where two gluons were radiated and these gluons would not necessarily be correlated in space.  It is therefore important to understand the effect of these kinematic cuts on the measured signal in various models for the production of the \ridge.

\section{Untriggered di-hadron correlations}

\begin{figure}
  \centering\resizebox{12cm}{!}{
  \includegraphics{/home/madscientistgirl/HarrisLab/Thesis/latex/previousresults/mikedaugheritypart1.eps}
}\\
  \centering\resizebox{12cm}{!}{
  \includegraphics{/home/madscientistgirl/HarrisLab/Thesis/latex/previousresults/mikedaugheritypart2.eps}
}\\

  \caption[Di-hadron correlations without \pT cuts in \Au collisions at \sNNtwohundred]
  {Di-hadron correlations without \pT cuts for different centralities in \Au collisions at \sNNtwohundred \cite{Daugherity:2008su}.  The z-axis is the probability for two particles to be correlated.
  \label{Mike}}
\end{figure}

Di-hadron correlations have also been done with no cuts on the momenta of either the trigger or the associated particles.  The motivation behind restricting the \pT of hadrons studied to high-\pT is to enhance the fraction of pairs coming from jets and reduce the background.  Samples of correlations without \pT cuts are shown in \Fref{Mike}.  When studies are done without cuts, all correlations between particles are evident, including those due to anisotropic flow, resonance decays, Hanbury-Brown-Twiss (HBT) correlations, Dalitz decays, conversion electrons, and jets.  Anisotropic flow leads to a cos(2\dphi) independent of \deta.  There is almost no contribution from anisotropic flow in the most peripheral collisions and the contribution is largest in the 19-28\% central events.  Dalitz decays, conversion electrons, and HBT correlations all lead to the sharp peak at \dphi=0, \deta=0.  The contributions from resonances is greatest near \dphi=0, \deta=0 and create structures that are dependent on the decay kinematics.  Resonance contributions were determined to be negligible for unidentified hadrons.  The contribution from jets contributes 2 peaks, one on the \ns and one on the \as.  Both of these peaks are evident in the most peripheral bin.   The whole correlation is fit to determine the relative contributions from these sources.  The \ns jet is the dominant feature near \dphi=0, \deta=0 in peripheral collisions.  This feature evolves from peripheral to central collisions into a peak which is much broader in pseudorapidity.  This feature may or may not be the same as the \ridge observed for high-\pT triggered correlations.

\section{Jet reconstruction}

\begin{figure}
  \centering\resizebox{8cm}{!}{
  \includegraphics{/home/madscientistgirl/HarrisLab/Thesis/latex/previousresults/ktjetsevil.eps}}\\

  \caption[Jet spectrum in central \Au collions at \sNNtwohundred]
  {Reconstructed jet spectrum in central \Au collisions at \sNNtwohundred as compared to the jet spectrum in \pp scaled by \Nbin \cite{jetrecoSevilHQ,Abelev:2007vt}.  Only statistical errors are shown for \Au data.  The systematic errors on the overall scale of the \pp spectrum are 50\%.
  \label{JetSpectra}}
\end{figure}

Full jet reconstruction allows determination of the parton energy and therefore would allow more careful measurements of partonic energy loss.  It also overcomes the uncertainty of whether or not a hadron comes from a jet, one of the most significant limitations of di-hadron correlations.  Full jet reconstruction was long thought to be impossible in heavy ion collisions at RHIC energies because the background is large and collisions were not expected to produce jets at a high enough energy to stand out above the background.  However, there have been improvements in the background subtraction method \cite{Cacciari:2008gn,Cacciari:2006sm} and the STAR collaboration has recently successfully reconstructed jets in heavy ion collisions \cite{jetrecoSevilWWND,jetrecoJoernHP,jetrecoSevilHP,jetrecoSevilHQ}.  \Fref{JetSpectra} shows the jet spectrum in central \Au collisions at \sNNtwohundred compared to the jet spectrum in \pp collisions scaled by the number of binary collisions.  If all of the energy of the jet is measured, the spectra in \Au and in \pp collisions scaled by \Nbin should be the same, as it appears to be.  This does not mean that the parton did not lose energy, only that this energy must appear somewhere because of energy conservation.  The systematic error on the \pp spectrum is 50\% \cite{Abelev:2007vt} and the systematic errors on the \Au spectrum are not yet determined.  The jets' shape in azimuth and pseudorapidity could still be modified and the fragmentation functions are not necessarily vacuum fragmentation functions.

Jet reconstruction and correlation studies are complementary.  If the shape is changed too dramatically, jet reconstruction algorithms may not recognize the modified jets.  Additionally, some studies of rare particles in jets are best suited to correlation studies due to statistics.

\section{Summary}

Studies of di-hadron correlations and nuclear modification factors at RHIC have demonstrated that jets are considerably modified by the medium.  This could be through partonic energy loss before fragmentation or through modification of the fragmentation functions in the medium.  These modifications have typically been classified as being on the \ns or the \as.  A note of caution is warranted since the \ns and the \as are defined only by their relation to the trigger particle and the only indication that the trigger particle might come from a jet is its momentum.

For an intermediate \pT trigger (roughly \trigrange{3.0}{6.0}), particles on the \as are observed at lower momentum than in \pp and \dAu collisions.  The \as is not observed at \ptassoc \roughly 2 \GeV (\Fref{IconicJetQuench}), but at lower momenta the \as is broadened and develops a double-peaked structure (\Fref{LowPtCorr}).  When the momenta of the trigger and associated particles are increased, the \as reappears (\Fref{PunchThrough}); this is interpreted as a greater bias towards unmodified jet fragmentation, perhaps through greater surface bias.  The \ridge is the dominant feature on the \ns for intermediate \pT triggers (\Fref{IconicRidgePlot}).  It persists out to \pttrig \roughly 8 \GeV, however, a greater percentage of the particles correlated with the trigger particle are in the \jlc (\Fref{RidgeVSPTTrigVSCent}).  As \ptassoc is increased, a greater percentage of the particles correlated with the trigger particle are in the \jlc (\Fref{RidgeVSJetVSPtAssoc}); at lower \ptassoc, measurements are overwhelmed by systematic errors and precise data are unavailable.  The features observed in correlation studies are momentum dependent, presumably reflecting differences in the relative importance of various mechanisms which produce correlations as a function of \pT.

The \jlc shows the \npart, \ptassoc, and \pttrig dependence that would be expected from vacuum fragmentation, although a more quantitative understanding of the \jlc is needed.  The \ridge is a new feature that appears to be made of particles from the bulk due to its similar composition (\Fref{RidgeParticleComposition}).  The spectra of particles in the \ridge are comparable to the inclusive spectra of particles, and the ratios of particles in the \ridge are similar to those in the bulk.  The \ridge is clearly correlated in space with a high-\pT particle, but that does not prove that the production of the two are related.  The full breadth of the data and an understanding of various models for the production of the \ridge are necessary to determine what this feature is, where it comes from, and what it can tell us about the QGP and relativistic heavy ion collisions.  The studies presented in this thesis will add to our understanding of the \jlc and the \ridge by investigating their dependence on collision energy and system size and their particle type dependence.